\section{Introduction}\label{sec:introduction}

Contact between the bodies of a multibody system is most often modelled at a point: a penalty
force that grows with the penetration depth, in the tradition of Hertz's theory
\cite{Hertz1882,HuntCrossley1975,Lankarani1990}, or a unilateral constraint with Coulomb friction
solved as a complementarity or a convex problem \cite{Stewart1996,Anitescu1997,Todorov2014}. Both
assume that the contact patch is small against the bodies. The reviews of contact-force models in
multibody dynamics describe the many variants of the first kind and their common ground in
non-conforming contact \cite{Machado2012,Skrinjar2018,Corral2021,Flores2022}.

Many contacts of practical interest do not satisfy that assumption. A flat part rests on a table,
a pin sits in a hole, a journal turns in a bearing, a stack of boxes leans, a gripper holds an
object between two pads. In these \emph{conforming} contacts the patch has an extent that
matters: the friction torque of a spinning disk, the tilt of an eccentric load, the angle at
which a tall body tips over a soft support or the force that holds a press fit all depend on how
the pressure is distributed over the patch. A point-contact model either does not represent
those effects or represents them through extra parameters, such as a torsional friction
coefficient with units of length, that the user must supply for each geometry.

Several families of models address the patch directly. Elastic-foundation models cover a body
with a bed of independent springs \cite{Johnson1985} and sum their forces over a discretized
surface; the polygonal contact model \cite{Hippmann2004} and Simbody's elastic foundation force
\cite{Sherman2011} are its established forms in multibody simulation. Volumetric models
\cite{Gonthier2005,Gonthier2007thesis} integrate the same foundation in closed form for a flat
patch and express the contact wrench through the volume, the centroid and the inertia tensor of
the overlap of the two undeformed bodies. Hasegawa and Sato \cite{Hasegawa2004} integrate a
distributed spring and damper over the exact intersection of convex polyhedra. Pressure-field
(hydroelastic) contact \cite{Elandt2019,Masterjohn2022} defines a pressure field inside each
body and a contact surface where the two fields are equal, and it is the patch model of the
Drake toolbox \cite{Drake}.

These models have been validated by their authors on a few cases each, and the published
comparisons of simulators do not fill the gap. The comparisons from robotics test solvers on
point contacts, by consistency under refinement, by physical criteria or against recorded
impacts \cite{Erez2015,Horak2019,Acosta2022,LeLidec2024}. The benchmark tradition of the
multibody community \cite{Gonzalez2006,Gonzalez2010,IFToMMBenchmark} is built on mechanisms with
ideal joints. Conformal contact has been studied where it matters most to that community, in
revolute joints with clearance, with finite-element and analytic references for the cylindrical
contact force \cite{Liu2007,Pereira2011,Pereira2015}, and comparisons of contact formulations
and of frictional-contact solvers exist \cite{Drumwright2011,Acary2018}. We did not find a set of
conforming-contact problems with analytic references that has been run across model families
and across engines.

This paper offers such a set and uses it. Its origin is the second author's doctoral thesis
\cite{Vilela2018}, which evaluated Gonthier's volumetric model on the exact intersection of two
triangle meshes inside one multibody code and tested it on three problems. The work reported
here re-implements that pipeline as a library that does not depend on the multibody engine,
adds a second law that integrates the same elastic foundation over the contact surface instead
of lumping it, and runs both, next to two baselines from the literature and the native contact
models of three engines, on a battery of problems with known answers.

The paper is a benchmark and a comparison; it does not propose a new contact model. The second
law, which we call \law{distributed}, belongs to the line that goes from Hasegawa and Sato
\cite{Hasegawa2004} through Gonthier \cite{Gonthier2005}, Hippmann's polygonal contact model
\cite{Hippmann2004} and Simbody's elastic foundation \cite{Sherman2011} to hydroelastic contact
\cite{Elandt2019}: the same Winkler foundation, here integrated exactly over the faces of the
polyhedral overlap. Section~\ref{sec:related} states what is taken from that line. The
contributions we claim are the following.
\begin{enumerate}
\item A battery of conforming- and patch-contact problems with analytic references, each
  classified by what its reference assumes (classes A, B and C), with a band of admissible
  pressure distributions for the references that assume one (Section~\ref{sec:benchmark}).
\item The findings, which are mechanisms: what a point contact cannot represent, what the
  lumping of Gonthier's model costs, and the failures of the two laws of this paper, each with
  a prediction that the runs match (the rolling resistance of the damped foundation, a
  friction-induced flutter with its stability condition, the stability bound of the explicit
  coupling, the compliance of the foundation) (Section~\ref{sec:results}).
\item A closed form for the Coulomb friction force and moment of a linear pressure on a
  triangle when the slip is a rigid motion (Section~\ref{sec:friction},
  Appendix~\ref{app:polar}). It is a short technical result: the closed form for uniform
  pressure on a polygon is known \cite{Jia1997}, and the linear term uses the same primitives.
  Earlier work on distributed contact interpolates the friction from the vertices or samples it
  at one point per face \cite{Hasegawa2004,ElandtThesis2022}. What the closed form gives is a
  friction wrench that does not depend on the triangulation of a flat face or on how a body is
  cut into pieces; element quadratures reach the same accuracy only on a refined, well-shaped
  mesh (Section~\ref{sec:ablation}).
\item An open-source, engine-independent implementation of both laws on the exact intersection
  of triangle meshes, with backends for Drake, Project Chrono and MuJoCo
  (Section~\ref{sec:methods}). ``Exact'' refers to the polyhedra given and to the normal
  wrench; the friction is exact for the linear equivalent of the pressure on each triangle.
\end{enumerate}

Section~\ref{sec:limitations} states the limits of this evidence in one place; the main ones are
that every reference is analytic and that the configurations of the native models are ours. The
constructions, the stability analysis, the native configurations with their results, the
complete tables and the statistical inference are in the supplementary material, cited as
Sections~S1 to~S5.

This work, from the implementation to the analysis and the text, was carried out by a
generative AI system, the model Claude Fable~5.2, listed as first author, under the direction and
supervision of the second author; the Declarations give the details.
